\section{Curvature Operator \texorpdfstring{$\Phi_7$}{Phi7}}
\label{sec:curvature}

This appendix provides the full mathematical development of the curvature
operator $\Phi_7$, corresponding to altitude $A7$ and domain $D_3$ in the
AxiomQuant Five Redux system. The operator acts on the spectral output
$s(x) = \Phi_5(x)$ and lifts it into curvature altitude via contraction with
the Riemann curvature tensor.

\subsection{Curvature Domain}

The curvature domain $D_3$ consists of all points $x \in \mathcal{M}$ equipped
with the Levi-Civita connection $\nabla$ and curvature tensor $R^i_{\ jkl}$.

Formally:
\[
D_3 = \{ x \in \mathcal{M} \mid R^i_{\ jkl}(x) \text{ is defined and smooth} \}.
\]

Altitude alignment:
\[
A7 \leftrightarrow D_3.
\]

\subsection{Riemann Curvature Tensor}

Recall the definition:
\[
R(X,Y)Z = \nabla_X \nabla_Y Z - \nabla_Y \nabla_X Z - \nabla_{[X,Y]} Z.
\]

In coordinates:
\[
R^i_{\ jkl}
= \partial_k \Gamma^i_{\ jl}
- \partial_l \Gamma^i_{\ jk}
+ \Gamma^i_{\ km}\Gamma^m_{\ jl}
- \Gamma^i_{\ lm}\Gamma^m_{\ jk}.
\]

Smoothness assumption:
\[
R^i_{\ jkl} \in C^\infty(\mathcal{M}).
\]

\subsection{Curvature Envelope}

Given the spectral signature
\[
s(x) = (\phi_1(x), \dots, \phi_m(x)),
\]
define the curvature envelope $c(x)$ by contracting $s(x)$ with the curvature
tensor.

Formally:
\[
c(x) = \Phi_7(s(x))
= \left( R^i_{\ jkl}(x) \, v^j v^k v^l \right)_{i=1}^{80},
\]
where $v$ is the spectral direction:
\[
v = \sum_{k=1}^m \phi_k(x) e_k.
\]

Thus:
\[
c_i(x) = R^i_{\ jkl}(x) \, v^j v^k v^l.
\]

This produces a curvature-modulated vector in $\mathbb{R}^{80}$.

\subsection{Operator Definition}

Formally:
\[
\Phi_7 : \mathbb{R}^m \to \mathbb{R}^{80},
\qquad
\Phi_7(s(x)) = c(x).
\]

Altitude alignment:
\[
\Phi_7 : A5 \to A7.
\]

Domain alignment:
\[
\Phi_7 : D_2 \to D_3.
\]

\subsection{Curvature Contraction Identity}

The contraction used by $\Phi_7$ satisfies:
\[
c_i = R^i_{\ jkl} v^j v^k v^l
= \langle R(\cdot, v)v, v \rangle_i.
\]

This identity is central to the geometric lift from spectral to curvature
altitude.

\subsection{Curvature Smoothness}

Since $R^i_{\ jkl}$ and $v^j$ are smooth:
\[
\Phi_7 \in C^\infty(\mathcal{M}, \mathbb{R}^{80}).
\]

\subsection{Curvature Stability}

For sufficiently small perturbations of the metric:
\[
\| \Delta g \| < \varepsilon
\quad\Longrightarrow\quad
\| \Delta R \| < C_1 \varepsilon,
\quad
\| \Delta v \| < C_2 \varepsilon,
\]
and thus:
\[
\| \Delta c(x) \| < C \varepsilon.
\]

\subsection{Curvature Reversibility}

On the curvature domain $D_3$, the map $\Phi_7$ is injective up to curvature
symmetry:
\[
\Phi_7(s(x)) = \Phi_7(s(x')) \Rightarrow x \sim x'.
\]

Higher-altitude operators resolve this equivalence.

\subsection{Curvature Monotonicity}

Altitude monotonicity:
\[
s(x) \le s(x') \Rightarrow \Phi_7(s(x)) \le \Phi_7(s(x')).
\]

This follows from the monotonicity of curvature under metric expansion.

\subsection{Operator Role}

The curvature operator provides the second geometric lift in the Redux
pipeline:
\[
s(x) \mapsto c(x),
\]
feeding directly into the holonomy operator $\Phi_8$.

This completes the formal development of the curvature operator.
